\scanpage{81}
\diagram{\begin{tikzcd}[ampersand replacement=\&]
X\arrow[d,"f"]\\ S
\end{tikzcd}}
\deleted{\unsure{propre}} plat de prés. \deleted{\unsure{pro}} finie

\[
\underline{\mathrm{Div}}_{X/S}\simeq
\underline{R}^{*}_{X/S}/\underline{O}^{*}_{X\deleted{\unsure{/S}}}
\xrightarrow{\frac{df}{f}}\underline{R}_{X/S}(\underline{\Omega}^1_{X/S})/\underline{\Omega}^1_{X/S}
\tag*{(*)}
\]
\deleted{\unsure{extension}}

d'où une extension [des couples $(D,\omega)$ tels que \deleted{\unsure{$dD$}} $d\log D\equiv\omega\;(\underline{\Omega}^1_{X/S})$]
\[
0\to\underline{\Omega}^1_{X/S}\to\underline E_{X/S}\to\underline{\mathrm{Div}}_{X/S}\to0
\]
\margin{\uline{N.B}. L'hom (*) est injectif si $S$ de car. nulle et $f$ normal \deleted{les fibres géométriques \unsure{normales}} (\unsure{loc.|donc} de prés. finie\ldots.)}

et par suite une suite exacte
\diagram{\begin{tikzcd}[ampersand replacement=\&,column sep=small]
0\arrow[r]\&f_*(\underline{\Omega}^1_{X/S})\arrow[r]\arrow[d,equal]\&f_*(\underline E_{X/S})\arrow[r]\&f_*(\underline{\mathrm{Div}}_{X/S})\arrow[r,"\partial"]\&R^1 f_*(\underline{\Omega}^1_{X/S})\\
\&\underline{\omega}_{X/S}
\end{tikzcd}}
\margin{à vérifier que c'est l'homomorphisme bien connu déduit de $\underline{\mathrm{Pic}}_{X/S}\to R^1 f_*(\underline{\Omega}^1_{X/S})$. c'est trivial}

ou encore
\diagram{\begin{tikzcd}[ampersand replacement=\&,column sep=small]
0\arrow[r]\&\underline{\omega}_{X/S}\arrow[r]\&f_*(\underline E_{X/S})\arrow[r]\&f_*(\underline{\mathrm{Div}}_{X/S})^{\unsure{w|W}}\arrow[r]\arrow[d,equal]\&0\\
\&\&\&\operatorname{Ker}(\partial\colon f_*(\underline{\mathrm{Div}}_{X/S})\to R^1 f_*(\underline{\Omega}^1_{X/S}))
\end{tikzcd}}

Considérons
\[
\underline P_{X/S}=f_*(\underline R^*_{X/S})/f_*(\underline O^*_X)\hookrightarrow f_*(\underline{\mathrm{Div}}_{X/S})\unsure{\leftarrow|-}
\]
On a, du moins si $X\to S$ est \uuline{quasi-projectif} \deleted{\unsure{et}}
\[
f_*(\underline{\mathrm{Div}}_{X/S})/\underline P_{X/S}\simeq\underline{\mathrm{Pic}}_{X/S}
\]
\deleted{d'autre part, \unsure{on a}}

et on pose
\[
\begin{aligned}
f_*(\underline{\mathrm{Div}}_{X/S})^w/\underline P_{X/S}
&=\underline{\mathrm{Pic}}^w_{X/S}\subset\underline{\mathrm{Pic}}_{X/S}\\
&=\operatorname{Ker}(\underline{\mathrm{Pic}}_{X/S}\to R^1 f_*(\underline{\Omega}^1_{X/S}))
\end{aligned}
\]

$*$ On \unsure{adjoint|définit} la flèche oblique
% archive sheet number: 80

